\begin{sang}{Bingospilleren}{FVJ (Finn Verner Jensen) - Mel: Manden på risten}
\spal 2
\begin{vers}
Jeg sidder hver torsdag med pladen, 
med palden til bingospil. 
Jeg håber at vinde til maden, 
lidt held er nødvendigt hertil. 
Min plade og jeg kan snart feste: 
Ti år har den været min. 
Jeg fik den i arv fra min bedste, 
hun arvede den fra sin. 
Den giver mig dejlige timer 
med spænding og ophidset blod, 
og ofte af gælde jeg hviner. 
Min palde er tro og god. 
Jeg glemmer bekymringer og sorger 
i selskab med dette ark. 
Når fat'nen er lykkeligt ovre, 
så får den et kys som tak. 
Så sig dog ikke, min verden er trist og banal, 
jeg finder lykke hos en plade med hele tal.
\end{vers}
\vfill
\begin{vers}
I tirsdags var jeg med min søster 
til bingo i Aabenaa, 
men ak, på sp fremmede kyser 
min plade ej bruges må. 
Jeg svigted' min plade i svaghed. 
Er der noget, man ikke skal, 
så er det at gi' sig i lag med 
en fremmed med lokkende tal. 
Der var 20 og 30 og 11 
og 66, 17 og 3. 
Jeg sad hele aft'nen og skælved', 
men præmierne udeble'. 
Imens gav jeg mig til at spille 
i tankerne med min eg'n. 
Jeg så til min skræk, at jeg ville 
ha' vundet en graveregn. 
Det sker ej mere. Nu har jeg truffet mit valg: 
Jeg svigter aldrig mere min plade med hele tal.
\end{vers}
\laps
\end{sang}